%% ch02.tex — Chapter 2: Data Integration: Heterogeneity, Quality, and Provenance
%% Part of "Data Engineering" textbook

\chapter{Data Integration: Heterogeneity, Quality, and Provenance}
\label{ch:integration}

\epigraph{\textit{``The hardest part of data science isn't machine learning or statistics.
It's getting data from different places to agree with each other.''}}{Anonymous data engineer}

\vspace{4pt}

\begin{chapterlearning}
After completing this chapter, you will be able to:
\begin{enumerate}
  \item Explain the data silo problem and articulate why integration
        complexity---not storage cost---is the primary barrier to
        enterprise analytics.
  \item Define the three fundamental operations of data integration:
        schema mapping, data exchange, and entity resolution.
  \item Distinguish among the four dimensions of data heterogeneity---
        structural, semantic, syntactic, and scale---and give a concrete
        example of each.
  \item Describe the six dimensions of data quality and identify
        appropriate remediation strategies for each.
  \item Explain what data provenance is, why it is required by regulations
        such as GDPR and HIPAA, and how lineage tracking systems capture it.
  \item Compare the ETL and ELT paradigms: their architectures, trade-offs,
        and appropriate use cases.
  \item Identify the right tool (Sqoop, Kafka, Airbyte, dbt) for each
        stage of a modern integration pipeline.
\end{enumerate}
\end{chapterlearning}

\bigskip

Chapter~\ref{ch:foundations} established the 3Vs as the defining
characteristics of big data: Volume, Velocity, and Variety. Volume and
Velocity receive the most attention in the popular press because they are
easy to quantify---petabytes of storage, thousands of events per second.
But \emph{Variety} is arguably the hardest engineering challenge. Data
arrives from dozens of independent sources, each with its own schema,
format, quality level, and update cadence. Before a single analytical
query can be answered, all of this heterogeneous data must be identified,
mapped, cleaned, and unified. This process is called \term{data
integration}, and it is the subject of this chapter.

A 2019 Forrester Research survey found that 73\% of enterprise data is
never analysed. The primary reason is not a shortage of storage or
processing power---it is integration complexity. Data that cannot be
found, joined, or trusted cannot be analysed. Data engineering at
production scale is, to a greater degree than is often acknowledged, the
art of making heterogeneous data coherent.

%%====================================================================
\section{The Data Silo Problem}
\label{sec:silos}
%%====================================================================

Large organisations accumulate data in \term{data silos}: independent
systems built by separate teams, at different times, for different
purposes, using different technology stacks. Each silo does its job
well in isolation, but no silo has visibility into the others. The
result is an organisation that is simultaneously data-rich and
insight-poor.

Consider a representative global retail bank. A customer who has held
an account for five years may appear as a distinct record in at least
seven separate systems: the core banking RDBMS (account balances,
transactions), the mobile application database (session events,
in-app behaviour), the loan origination system (credit applications,
repayment history), the customer relationship management platform
(support tickets, call transcripts), the anti-money-laundering engine
(transaction flags, compliance cases), the marketing automation
platform (campaign responses, email open rates), and the regulatory
reporting database (filings with supervisory authorities). In each
system, the customer may be identified by a different key---an internal
account number, a mobile device identifier, a CRM contact ID, a
government-issued reference number---and the same attribute (the
customer's full name, for example) may be spelled differently, encoded
differently, or stored at different granularities.

The consequence is not merely an engineering inconvenience. It is a
fundamental barrier to analytical insight: a data scientist wishing to
understand the relationship between a customer's loan repayment
behaviour and their mobile application engagement cannot combine
these datasets without first resolving which records in the loan system
correspond to which records in the mobile database---a non-trivial
operation even for a single bank with thousands of customer
records, let alone for a global institution with hundreds of millions.

\begin{keyinsight}{Integration is the Critical Path}
In most data engineering projects, the time spent on data integration---
schema mapping, entity resolution, quality cleaning, and pipeline
construction---exceeds the time spent on the analytical computation
itself. Studies of data science workflows consistently find that
engineers spend 60--80\% of project time on data preparation rather
than on modelling or analysis. Integration quality directly determines
analytical quality: a machine learning model trained on incorrectly
joined data will encode the integration error into its predictions.
\end{keyinsight}

%%====================================================================
\section{A Formal Framework for Data Integration}
\label{sec:integration-framework}
%%====================================================================

The academic study of data integration has a rich history spanning
decades of database research. Halevy, Rajaraman, and Ordille's landmark
2006 survey---titled ``Data Integration: The Teenage Years''---synthesised
thirty years of work and established the formal vocabulary that the
field still uses today \parencite{Halevy2006}.

\begin{defbox}{Data Integration}
Given a set of heterogeneous, autonomous data sources
$S_1, S_2, \ldots, S_n$, each with its own schema $\Sigma_i$ and
content $D_i$, \textbf{data integration} is the problem of providing
a unified query interface over a \textbf{mediated (global) schema}
$\Sigma_G$ that gives users a single coherent view of all data,
regardless of which source it originated from.
\end{defbox}

This definition makes three important properties explicit. First, data
sources are \emph{heterogeneous}: they differ in schema, format,
vocabulary, and update cadence. Second, they are \emph{autonomous}: the
integration layer has no authority to modify the source systems---it can
only read from them. Third, the goal is a \emph{unified query interface}:
users should be able to ask a single query against the global schema
without knowing which source holds the relevant data or in what format.

Achieving this goal requires three core operations, each of which is a
research area in its own right.

\subsection{Schema Mapping}
\label{subsec:schema-mapping}

\term{Schema mapping} is the process of defining a formal correspondence
between the schema of a source system and the global mediated schema.
A schema mapping $\mathcal{M}_{i \to G}$ specifies, for each attribute
in the global schema, which attribute in source $S_i$ provides the
corresponding data, and what transformation (if any) is required to
reconcile differences in naming, data type, units, or encoding.

Schema mappings are expressed as rules. For example:
\begin{equation}
\texttt{GlobalSchema.customer\_id} \leftarrow \texttt{CRM.contact\_id}
\quad [\text{identity}]
\end{equation}
\begin{equation}
\texttt{GlobalSchema.full\_name} \leftarrow
\texttt{uppercase}(\texttt{ATM.CUST\_FNAME}
+ \texttt{" "} + \texttt{ATM.CUST\_LNAME})
\end{equation}

In practice, deriving accurate schema mappings is one of the most
labour-intensive tasks in data engineering. Automated \term{schema
matching} tools---which use string similarity, structural similarity,
and machine learning to propose candidate mappings---can reduce this
effort significantly, but they always require human validation before
production deployment.

\subsection{Data Exchange}
\label{subsec:data-exchange}

\term{Data exchange} is the process of \emph{materialising} data from
source schemas into the global schema. It is the implementation of a
schema mapping: given source data $D_i$ and a mapping
$\mathcal{M}_{i \to G}$, produce a dataset $D_G$ that conforms to the
global schema $\Sigma_G$. Data exchange corresponds to the
``Transform'' step in an ETL pipeline (discussed in
Section~\ref{sec:etl-elt}).

\subsection{Entity Resolution}
\label{subsec:entity-resolution}

\term{Entity resolution} (also called \term{record linkage},
\term{deduplication}, or \term{identity resolution}) is the process of
determining whether two records from different sources (or from the same
source) refer to the same real-world entity. It is necessary when
different systems use different identifiers for the same object and when
direct key-based joins are therefore impossible.

Entity resolution is examined in detail in Section~\ref{sec:er}.

\begin{researchspotlight}{Halevy, Rajaraman \& Ordille (2006) --- Data Integration: The Teenage Years}
\textbf{Citation:} Halevy, A., Rajaraman, A., \& Ordille, J.\ (2006).
Data Integration: The Teenage Years. \textit{Proceedings of the 32nd
VLDB Conference}, pp.\ 9--16. \parencite{Halevy2006}

\smallskip\noindent
\textbf{Key insight:} The paper's title is a pointed metaphor: after
three decades of research, data integration remained ``in its teenage
years''---having demonstrated proof of concept and produced fundamental
theory, but not yet reaching the robust, deployable maturity that
industry needed.

\smallskip\noindent
\textbf{Technical contributions:}
\begin{itemize}
  \item Provided a comprehensive taxonomy of the data integration
        problem space, distinguishing virtual integration (query
        rewriting over live sources) from materialised integration
        (ETL into a data warehouse).
  \item Formalised the distinction between \emph{local-as-view} (LAV)
        and \emph{global-as-view} (GAV) approaches to schema mapping---
        still the canonical framework for integration system design.
  \item Identified entity resolution as the hardest open problem in
        data integration, predicting (correctly) that it would require
        probabilistic and machine learning approaches at scale.
\end{itemize}

\smallskip\noindent
\textbf{Impact today:} The paper is required reading in database
courses at MIT, Stanford, and CMU. Its taxonomy of integration
challenges directly influenced the design of modern enterprise data
integration platforms, including MuleSoft, Informatica, and Apache
Atlas. The LAV/GAV distinction remains the theoretical foundation of
modern data virtualisation products.
\end{researchspotlight}

%%====================================================================
\section{The Four Dimensions of Data Heterogeneity}
\label{sec:heterogeneity}
%%====================================================================

Heterogeneity---the fundamental barrier to integration---manifests along
four distinct dimensions. Understanding all four is essential because a
pipeline that addresses only one or two will produce subtly incorrect
results, often without any obvious indication of failure.

\begin{figure}[h!t]
\centering
\begin{tikzpicture}[
    quad/.style={rectangle, rounded corners=5pt,
                 minimum width=4.8cm, minimum height=1.8cm,
                 text width=4.5cm, align=center,
                 font=\small\bfseries},
    sub/.style={font=\scriptsize\normalfont},
    arr/.style={-{Stealth[length=4pt]}, thick, DEBlue!40}
  ]
  \node[quad, fill=DEBlue,    text=white] (A) at (0,   2.2)
        {Structural\\[3pt]{\small\normalfont Different table/column structures}};
  \node[quad, fill=DEAccent,  text=white] (B) at (5.5, 2.2)
        {Semantic\\[3pt]{\small\normalfont Same word, different meaning}};
  \node[quad, fill=DEGreen,   text=white] (C) at (0,   0)
        {Syntactic\\[3pt]{\small\normalfont Same data, different encoding}};
  \node[quad, fill=DEPurple,  text=white] (D) at (5.5, 0)
        {Scale \& Granularity\\[3pt]{\small\normalfont Different units or resolutions}};

  % Centre hub
  \node[circle, fill=DEGray!20, draw=DEGray!50, minimum size=1.1cm,
        font=\footnotesize\bfseries\sffamily, align=center, text=DEBlue]
        (hub) at (2.75, 1.1) {Data\\Hetero-\\geneity};

  \draw[arr] (A.east |- hub) -- (hub);
  \draw[arr] (B.west |- hub) -- (hub);
  \draw[arr] (C.east |- hub) -- (hub);
  \draw[arr] (D.west |- hub) -- (hub);
\end{tikzpicture}
\caption{The four dimensions of data heterogeneity. A production
integration pipeline must address all four simultaneously.}
\label{fig:heterogeneity-types}
\end{figure}

%--------------------------------------------------------------------
\subsection{Structural Heterogeneity}
\label{subsec:structural-hetero}
%--------------------------------------------------------------------

\term{Structural heterogeneity} arises when two data sources represent
the same real-world concept using different table organisations, column
names, or normalisation levels. It is the most immediately obvious form
of heterogeneity and the first that integration engineers typically
address.

\textbf{Naming conflicts} occur when the same attribute has different
names in different systems. A customer identifier might be called
\texttt{cust\_id} in the core banking system, \texttt{customer\_number}
in the loan system, and \texttt{contact\_id} in the CRM. Without a
mapping, a JOIN across these tables produces no results.

\textbf{Structural conflicts} occur when the same data is stored at
different levels of normalisation. One system may store a customer's
address as a single string (\texttt{"14 Oxford Street, London, W1D 1AU"}),
while another uses four separate columns (\texttt{street\_line\_1},
\texttt{city}, \texttt{region}, \texttt{postcode}). Converting between
these representations requires both parsing and composing---two
operations that are error-prone when addresses do not conform to
expected patterns.

\begin{table}[h!t]
\centering
\caption{Structural heterogeneity: a customer record in three systems}
\label{tab:structural-hetero}
\renewcommand{\arraystretch}{1.3}
\begin{tabular}{L{3.0cm} L{2.8cm} L{2.6cm} L{2.4cm}}
\toprule
\textbf{Concept} & \textbf{Core Banking} & \textbf{CRM} & \textbf{ATM Logs} \\
\midrule
Customer ID      & \texttt{cust\_id}        & \texttt{contact\_id}  & \texttt{CARD\_HOLDER} \\
First name       & \texttt{first\_name}     & \texttt{fname}        & \texttt{FN} \\
Date of birth    & \texttt{dob} (DATE)      & \texttt{birth\_date} (VARCHAR) & absent \\
Address          & 4-column normalised      & Single string         & absent \\
Account open date & \texttt{open\_date}    & \texttt{create\_ts}   & absent \\
\bottomrule
\end{tabular}
\end{table}

%--------------------------------------------------------------------
\subsection{Semantic Heterogeneity}
\label{subsec:semantic-hetero}
%--------------------------------------------------------------------

\term{Semantic heterogeneity}---the most dangerous form---arises when
two data sources use the same name to mean different things, or
different names to mean the same thing. A schema mapping that resolves
structural conflicts (renaming columns) will silently produce incorrect
results if the underlying semantic conflict is not also resolved.

Classic examples are pervasive:
\begin{itemize}
  \item The column \texttt{price} in the inventory system records the
        wholesale cost; the same column in the e-commerce system
        records the retail price. Joining these tables without
        disambiguation produces a dataset where ``price'' has a
        different meaning for every row.
  \item The field \texttt{date} in an order management system could
        mean the order date, the shipping date, or the invoice date,
        depending on which team created the table. All three share
        the same column name in different tables.
  \item The concept of a ``customer'' means a natural person in the
        retail banking system, but can also mean a legal entity
        (company) in the corporate banking system. A JOIN on
        \texttt{customer\_id} would mix individuals and corporations.
\end{itemize}

Semantic heterogeneity cannot be detected by examining schemas alone---
it requires domain knowledge, documentation review, and often direct
consultation with the teams that created each system. This is why data
dictionaries and business glossaries are not administrative overhead but
engineering necessities.

%--------------------------------------------------------------------
\subsection{Syntactic Heterogeneity}
\label{subsec:syntactic-hetero}
%--------------------------------------------------------------------

\term{Syntactic heterogeneity} arises when the same data is encoded in
different formats, character encodings, or data types across sources.
The value and meaning are equivalent, but the byte-level representation
differs. Syntactic heterogeneity must be resolved before any comparison
or JOIN can succeed.

Representative examples include:
\begin{itemize}
  \item Dates stored as \texttt{"2024-03-15"} (ISO 8601, universally
        sortable) in one system and as \texttt{"15/03/2024"} (UK format)
        or \texttt{"03-15-2024"} (US format) in another. A lexicographic
        sort on the date column produces incorrect ordering when formats
        are mixed.
  \item Boolean values represented as \texttt{true/false} (JSON),
        \texttt{1/0} (integer), \texttt{"Y"/"N"} (legacy COBOL systems),
        or \texttt{"yes"/"no"} (human-entered data).
  \item Currency amounts stored as 64-bit floating-point numbers in one
        system and as decimal strings (e.g., \texttt{"1,234.56"} or
        \texttt{"1.234,56"} for European locale) in another.
  \item Character encodings: legacy mainframe systems frequently use
        EBCDIC encoding; modern web applications use UTF-8; some regional
        systems use Windows-1252 or ISO-8859 variants.
\end{itemize}

%--------------------------------------------------------------------
\subsection{Scale and Granularity Heterogeneity}
\label{subsec:scale-hetero}
%--------------------------------------------------------------------

\term{Scale and granularity heterogeneity} arises when different sources
measure the same quantity at different units, resolutions, or temporal
granularities. Even when structural, semantic, and syntactic conflicts
have been resolved, a mismatch in granularity can invalidate
comparisons.

A revenue report that joins monthly aggregates from one system with
daily transaction records from another must aggregate the latter before
joining; failing to do so produces a systematic distortion where
each monthly figure appears 30 times. Similarly, a climate analytics
system integrating temperature readings from stations reporting in
Celsius and others in Fahrenheit must apply a unit conversion before
any cross-station analysis; a system integrating GPS coordinates at
metre precision with satellite imagery at 10-metre resolution must
account for the resolution mismatch when attributing imagery to
ground-truth measurements.

\begin{techdebt}{Integration Failures Are Silent}
Of the four heterogeneity dimensions, structural and syntactic failures
are the easiest to detect: they typically produce explicit errors
(column not found, type mismatch, parse failure). Semantic and
granularity failures are far more dangerous because they produce
\emph{no error}---the pipeline succeeds, data flows, and the output
looks superficially correct. A data warehouse that has silently mixed
wholesale prices with retail prices will answer every price-related
query without complaint, producing reports that are confidently and
systematically wrong. Detecting semantic errors requires domain
knowledge and statistical profiling of the integrated output, not
just schema validation.
\end{techdebt}

%%====================================================================
\section{Entity Resolution}
\label{sec:er}
%%====================================================================

Even after schema mappings have been defined and all four heterogeneity
dimensions addressed, a fundamental problem remains: how do we know that
record $r_1$ from source $S_1$ and record $r_2$ from source $S_2$ refer
to the same real-world entity? When a common key exists (both systems
use a social security number or a product EAN barcode), the answer is
trivial: join on the key. When no common key exists---the common
case---the problem is \term{entity resolution}.

\begin{defbox}{Entity Resolution}
Given two sets of records $R_1$ and $R_2$ from sources $S_1$ and $S_2$
respectively, \textbf{entity resolution} is the task of identifying all
pairs $(r_1, r_2) \in R_1 \times R_2$ that refer to the same
real-world entity, using only the attribute values of the records (and
no common key).
\end{defbox}

The naive approach---comparing every record in $R_1$ against every
record in $R_2$---has complexity $O(|R_1| \times |R_2|)$. For datasets
with millions of records, this is computationally infeasible. A dataset
of 10 million customer records from two sources would require 100
trillion comparisons---at a rate of one billion comparisons per second,
that is over 27 hours. At big data scale, the quadratic cost is
prohibitive.

%--------------------------------------------------------------------
\subsection{Blocking: Reducing the Candidate Space}
\label{subsec:blocking}
%--------------------------------------------------------------------

\term{Blocking} is the standard technique for making entity resolution
computationally tractable. Rather than comparing all possible pairs,
blocking partitions records into \term{blocks} based on a simple
attribute (or combination of attributes) that is likely to match for
true duplicates. Only pairs within the same block are compared. A common
blocking key is the first three characters of the last name combined
with the year of birth; two records can only be a match if they share
the same block key.

A well-designed blocking strategy must balance two competing
objectives. \emph{Recall} requires that almost all true duplicates fall
in the same block---if a blocking key is too aggressive, it splits true
duplicates into different blocks and they are never compared.
\emph{Efficiency} requires that blocks be small enough to make
comparison tractable---if a blocking key is too lenient, blocks are
large and the quadratic comparison cost returns. In practice, multiple
blocking keys are used in parallel, and a pair is a candidate if it
shares at least one blocking key.

%--------------------------------------------------------------------
\subsection{Similarity Functions}
\label{subsec:similarity}
%--------------------------------------------------------------------

Within each block, records are compared using \term{similarity
functions} that assign a score $\sigma(r_1, r_2) \in [0, 1]$ to each
candidate pair. Two commonly used functions are:

\textbf{Jaccard similarity} operates on sets of tokens. Given two
strings $s_1$ and $s_2$ tokenised into sets of words or character
n-grams $T_1$ and $T_2$:
\begin{equation}
J(T_1, T_2) = \frac{|T_1 \cap T_2|}{|T_1 \cup T_2|}
\label{eq:jaccard}
\end{equation}
Jaccard similarity is robust to word reordering and partial matches,
making it well-suited for company names and addresses.

\textbf{Levenshtein (edit) distance} measures the minimum number of
single-character insertions, deletions, or substitutions needed to
transform one string into another. Normalised edit similarity is:
\begin{equation}
\sigma_\text{edit}(s_1, s_2) = 1 -
  \frac{d_\text{edit}(s_1, s_2)}{\max(|s_1|, |s_2|)}
\label{eq:edit-sim}
\end{equation}
Edit distance handles typos, transpositions, and abbreviated names
(``Corp.'' vs ``Corporation'').

A \term{composite similarity score} combines multiple attribute
similarities with weights reflecting their discriminative power:
\begin{equation}
\sigma_\text{composite}(r_1, r_2) =
  w_\text{name} \cdot \sigma(r_1.\text{name},\, r_2.\text{name})
  + w_\text{addr} \cdot \sigma(r_1.\text{addr},\, r_2.\text{addr})
  + w_\text{dob}  \cdot \sigma(r_1.\text{dob},\, r_2.\text{dob})
\label{eq:composite-sim}
\end{equation}
where $w_\text{name} + w_\text{addr} + w_\text{dob} = 1$.

A pair $(r_1, r_2)$ is declared a match if
$\sigma_\text{composite}(r_1, r_2) \geq \theta$ for a threshold $\theta$
chosen to balance precision and recall for the specific application.

%%====================================================================
\section{Data Quality: Dimensions and Challenges}
\label{sec:data-quality}
%%====================================================================

Data integration can perfectly resolve structural, semantic, syntactic,
and granularity heterogeneity, correctly link every entity, and still
produce a dataset that is analytically worthless---if the underlying
data is of poor quality. The \term{Veracity} dimension of the 5Vs
framework (Chapter~\ref{ch:foundations}) captures this concern. Data
quality problems are not exceptions; they are ubiquitous in production
systems, and big data engineering requires systematic processes for
detecting, quantifying, and remediating them.

The field has converged on six dimensions of data quality, originally
formalised by Wang and Strong (1996) in their empirical study of
information consumers' quality requirements.

\begin{table}[h!t]
\centering
\caption{The six dimensions of data quality with engineering implications}
\label{tab:dq-dimensions}
\renewcommand{\arraystretch}{1.45}
\begin{tabular}{L{2.2cm} L{3.6cm} L{4.6cm}}
\toprule
\textbf{Dimension} & \textbf{Definition} & \textbf{Detection and remediation} \\
\midrule
Accuracy      & Values correctly represent the real-world entity or event &
                Statistical profiling; external reference datasets; domain
                range constraints \\
Completeness  & All required values are present; no critical nulls &
                NULL ratio checks; conditional completeness rules
                (\texttt{if type=``corporate'' then tax\_id must be non-null}) \\
Consistency   & Values are not contradictory within or across sources &
                Cross-source join validation; referential integrity checks;
                temporal monotonicity checks \\
Timeliness    & Data is current enough for its intended use &
                Staleness metrics; data freshness SLAs; arrival-time audits \\
Validity      & Values conform to defined formats, ranges, and vocabularies &
                Regular expression validation; enumeration checks;
                numeric range constraints \\
Uniqueness    & Records are not duplicated (within a source or across sources) &
                Duplicate detection; entity resolution (Section~\ref{sec:er}) \\
\bottomrule
\end{tabular}
\end{table}

%--------------------------------------------------------------------
\subsection{Missing Values}
\label{subsec:missing-values}
%--------------------------------------------------------------------

Missing values are among the most common data quality problems in large
datasets. They arise from several distinct mechanisms, each with
different analytical implications.

\textbf{Missing Completely at Random} (MCAR): the probability that a
value is missing is independent of both the observed and missing values.
A sensor that fails with a fixed, context-independent probability
produces MCAR missingness. Data imputation strategies---such as mean or
median imputation---are statistically valid under MCAR.

\textbf{Missing at Random} (MAR): the probability that a value is
missing depends on observed values but not on the missing value itself.
For example, older customers may be less likely to provide an email
address, so the \texttt{email} field is missing more often for older
cohorts. This is MAR because age (observed) predicts missingness, but
the missing value itself (specific email) does not.

\textbf{Missing Not at Random} (MNAR): the most problematic case. The
probability that a value is missing depends on the value itself.
Survey data where respondents with very low incomes decline to answer
the income question is MNAR---the missingness is informative about
the missing value. Imputing MNAR data with mean imputation introduces
systematic bias.

In big data engineering, the immediate engineering response to missing
values is typically a combination of: (1) quantifying the NULL rate
per column in the data quality dashboard, (2) identifying whether
the missing values are upstream (the source system never had the
value) or midstream (the ETL pipeline dropped them), and (3) applying
domain-appropriate imputation, filtering, or flagging. Feeding a machine
learning model records with missing feature values without appropriate
handling is one of the most common sources of production ML failure.

%--------------------------------------------------------------------
\subsection{Noise, Outliers, and Inconsistency}
\label{subsec:noise-outliers}
%--------------------------------------------------------------------

\textbf{Noise} refers to random errors in data values---measurement
errors, transmission errors, or OCR mistakes. A temperature sensor that
occasionally reads --999.0$^{\circ}$C due to a firmware bug introduces
noise; a mobile payment amount that loses a decimal point and appears
as 10,000 times the correct value is noise. At scale, even a 0.01\%
noise rate in a 10-billion-row dataset produces one million corrupted
records---enough to systematically distort any aggregate computed over
that column.

\textbf{Outliers} may be noise, or they may be genuine rare events.
A single transaction of \$50 million in a retail payment network is
an outlier; it might represent fraud (noise), or it might represent a
legitimate large commercial payment. In data engineering, the
distinction between noise and genuine outliers must be resolved using
domain knowledge and external validation, not purely statistical
criteria.

\textbf{Inconsistency} arises when the same real-world fact is
represented differently at different points in time or in different
parts of the same system. A customer record that shows the customer's
home city as ``New York'' in the CRM but ``NYC'' in the billing system
is inconsistent. A product price that changes in the catalogue but is
not updated in the recommendation engine creates a temporal
inconsistency. Inconsistency is particularly dangerous in slowly
changing dimension tables (a concept explored in the data warehousing
literature) where historical accuracy must be maintained alongside
current values.

%%====================================================================
\section{Data Provenance and Lineage}
\label{sec:provenance}
%%====================================================================

\term{Data provenance} refers to the documented record of a data
artifact's origin, transformation history, and movement through a
system---the complete answer to the question: ``Where did this value
come from, and what was done to it?'' \term{Data lineage} is the
operationalised form of provenance: the graph of dependencies linking
each output dataset to its input datasets and the transformations
applied to produce it.

Provenance is not merely a technical nicety; it is increasingly a
legal requirement. The European Union's General Data Protection
Regulation (GDPR, 2018) requires organisations to be able to demonstrate
the legal basis for processing personal data, trace where personal data
flows within and across systems, and respond to subject access requests
that require identifying all locations where an individual's data is
held. The United States' Health Insurance Portability and
Accountability Act (HIPAA) requires healthcare organisations to maintain
audit trails documenting every access to and movement of protected
health information (PHI). Failure to demonstrate adequate provenance
tracking has resulted in regulatory fines exceeding \$10 million in
individual GDPR enforcement actions.

Beyond regulatory compliance, provenance serves crucial operational
purposes. When an analytical model produces an unexpected result, data
lineage allows engineers to trace the result back to the exact source
records and transformations that produced it---identifying whether the
anomaly originates in a source system, an ETL transformation, or the
analytical model itself. In production data pipelines, this capability
can reduce the mean time to detect and resolve data quality incidents
from days to hours.

\begin{keyinsight}{Lineage as a First-Class Engineering Artifact}
Modern data engineering treats lineage as a first-class artifact, not
an afterthought. Tools such as Apache Atlas (metadata management),
OpenLineage (open specification for lineage events), and dbt's built-in
lineage graph capture transformations automatically as pipelines run.
A comprehensive lineage graph answers: which tables feed this model,
which source records produced this output row, which transformation
version was active when this job ran, and who approved the pipeline
change that introduced this transformation. Without this capability,
debugging a corrupted KPI dashboard may require days of manual source
tracing across dozens of systems.
\end{keyinsight}

%%====================================================================
\section{ETL vs.\ ELT: Two Philosophies of Integration}
\label{sec:etl-elt}
%%====================================================================

The practical execution of data integration---moving, transforming, and
loading data from source systems into analytical targets---is governed
by two contrasting paradigms. The choice between them is one of the most
consequential architectural decisions in data engineering, because it
determines where and when data is transformed, how flexible the system
is to evolving requirements, and what trade-offs are made between data
freshness, quality, and engineering complexity.

%--------------------------------------------------------------------
\subsection{ETL: Extract, Transform, Load}
\label{subsec:etl}
%--------------------------------------------------------------------

\term{ETL} is the traditional paradigm, dominant in data warehousing
from the 1980s through the 2010s. It operates in three sequential
phases:

\textbf{Extract}: Data is read from source systems---typically RDBMS
tables, file exports, or API responses---and written to a staging area.
The extraction is designed to be minimally invasive, typically using
incremental extraction (change data capture, CDC) rather than full table
scans, to avoid overloading source systems.

\textbf{Transform}: In the staging area, the extracted data is
subjected to all required cleaning, mapping, enrichment, and
aggregation operations before it is written to the target. The target
(typically a data warehouse) enforces a strict schema; data that does
not conform to the schema is rejected. This means that the transformation
step must be comprehensive and correct before any data reaches the
warehouse.

\textbf{Load}: The transformed, validated data is loaded into the
target data warehouse in bulk. Only data that passes all quality and
schema checks is written to the target.

The primary advantage of ETL is data quality at the target: by the time
data reaches the warehouse, it has been fully validated and conforms
precisely to the intended schema. The primary disadvantage is
rigidity: adding a new source or changing a transformation requires
modifying the ETL pipeline, which may require schema changes in the
warehouse, regression testing of all downstream queries, and potentially
a full historical reload.

%--------------------------------------------------------------------
\subsection{ELT: Extract, Load, Transform}
\label{subsec:elt}
%--------------------------------------------------------------------

\term{ELT} inverts the order: raw data is extracted from sources and
loaded immediately into a data lake (typically cloud object storage such
as Amazon S3, Google Cloud Storage, or Azure Data Lake Storage) in its
native format, without transformation. Transformation is deferred until
analytical consumption time, when a compute engine such as Apache Spark
or a transformation tool such as dbt applies the necessary logic
on-demand.

ELT has become the dominant paradigm for modern data engineering for
several reasons. First, cloud object storage is extraordinarily
inexpensive (approximately \$0.023 per GB per month for S3 Standard as
of 2024), making it economically feasible to store raw, untransformed
data at any scale. Second, decoupling extraction from transformation
enables \emph{schema evolution}: when a source system changes its schema,
the raw data in the lake is unaffected, and only the transformation
logic needs to be updated. Third, deferred transformation supports
multiple downstream use cases with different transformation requirements
without requiring the ETL pipeline to anticipate every future analytical
need.

\begin{figure}[h!t]
\centering
\begin{tikzpicture}[
    box/.style={rectangle, rounded corners=3pt, minimum width=2.4cm,
                minimum height=0.8cm, font=\small\bfseries\sffamily,
                align=center},
    pipe/.style={-{Stealth[length=5pt]}, very thick},
    lbl/.style={font=\scriptsize\sffamily, align=center}
  ]
  %--- ETL (top row) ---
  \node[font=\small\bfseries\sffamily\color{DEBlue}] at (0.0, 3.8) {ETL};
  \node[box, fill=DEBlue!20,   draw=DEBlue!60]  (e1) at (0.0, 3.0) {Sources};
  \node[box, fill=DEAccent!30, draw=DEAccent!70](t1) at (3.0, 3.0) {Staging \&\\Transform};
  \node[box, fill=DEGreen!30,  draw=DEGreen!70] (l1) at (6.0, 3.0) {Data\\Warehouse};
  \draw[pipe, DEBlue!70]   (e1) -- (t1) node[midway,above,lbl]{Extract};
  \draw[pipe, DEAccent!80] (t1) -- (l1) node[midway,above,lbl]{Load clean};

  %--- ELT (bottom row) ---
  \node[font=\small\bfseries\sffamily\color{DEBlue}] at (0.0, 1.8) {ELT};
  \node[box, fill=DEBlue!20,   draw=DEBlue!60]  (e2) at (0.0, 1.0) {Sources};
  \node[box, fill=DEGreen!30,  draw=DEGreen!70] (l2) at (3.0, 1.0) {Data Lake\\(raw)};
  \node[box, fill=DEAccent!30, draw=DEAccent!70](t2) at (6.0, 1.0) {Transform\\on query};
  \draw[pipe, DEBlue!70]   (e2) -- (l2) node[midway,above,lbl]{Extract + Load};
  \draw[pipe, DEGreen!80]  (l2) -- (t2) node[midway,above,lbl]{Transform};

  % Quality indicator
  \node[lbl, text=DEGreen!70!black] at (6.0, 2.4) {High quality\\strict schema};
  \node[lbl, text=DEAccent!80!black] at (6.0, 0.2) {Flexible\\schema-on-read};
\end{tikzpicture}
\caption{ETL vs.\ ELT: the order of operations determines where
and when data quality constraints are enforced.}
\label{fig:etl-elt}
\end{figure}

\begin{table}[h!t]
\centering
\caption{ETL vs.\ ELT: a systematic comparison}
\label{tab:etl-elt}
\renewcommand{\arraystretch}{1.4}
\begin{tabular}{L{2.6cm} L{3.4cm} L{4.0cm}}
\toprule
\textbf{Attribute}         & \textbf{ETL}                    & \textbf{ELT} \\
\midrule
Transformation timing      & Before loading (staging area)   & After loading (at query time) \\
Target schema              & Predefined, strict (schema-on-write) & Flexible (schema-on-read) \\
Data quality at target     & High (validated before load)    & Variable (raw data preserved) \\
Historical replay          & Expensive (re-run ETL jobs)     & Easy (re-transform from raw) \\
Schema evolution           & Difficult (pipeline rebuild)    & Easy (update transformation only) \\
Raw data preservation      & No (transformed data only)      & Yes (raw always available) \\
Storage cost               & Lower (transformed, compressed) & Higher (raw data at full volume) \\
Latency                    & Higher (transform-then-load)    & Lower (load immediately) \\
Tooling                    & Informatica, Talend, SSIS       & Spark, dbt, Airbyte, Fivetran \\
\textbf{Best suited for}   & Regulated industries, strict SLAs & Data lakes, ML feature stores, exploratory analytics \\
\bottomrule
\end{tabular}
\end{table}

\begin{caution}{ELT Does Not Mean Skip Transformation}
A common misunderstanding is that ELT eliminates the need for data
transformation. It does not---it defers transformation. The raw data in
a data lake is typically not directly queryable or usable for analytics;
it requires the same cleaning, normalisation, and enrichment that ETL
performs, just applied at query time rather than ingestion time. The
\emph{dbt} (data build tool) framework has become the standard for
managing this deferred transformation layer, implementing transformation
logic as versioned, tested SQL models with lineage tracking. ELT without
a robust transformation and testing layer produces a ``data swamp''---
a lake of raw data that nobody trusts or can efficiently query.
\end{caution}

%%====================================================================
\section{The Modern Data Integration Stack}
\label{sec:integration-stack}
%%====================================================================

Modern data integration pipelines combine tools that address different
stages and dimensions of the integration problem. Understanding the
ecosystem---which tool solves which problem---is essential for designing
production systems.

\begin{table}[h!t]
\centering
\caption{The modern data integration tool ecosystem}
\label{tab:integration-tools}
\renewcommand{\arraystretch}{1.4}
\begin{tabular}{L{2.2cm} L{2.0cm} L{2.0cm} L{4.8cm}}
\toprule
\textbf{Tool}          & \textbf{Direction}   & \textbf{Pattern} & \textbf{Primary use case} \\
\midrule
Apache Sqoop           & RDBMS $\leftrightarrow$ HDFS & Batch    & Bulk migration of relational tables into Hadoop \\
Apache Flume           & Log sources $\to$ HDFS & Batch/micro-batch & Web server log aggregation \\
Apache Kafka           & Any $\to$ Any         & Stream    & Real-time event bus; decoupled source-to-sink \\
Airbyte / Fivetran     & SaaS APIs $\to$ Data lake & Batch/incremental & No-code connector platform for 200+ sources \\
Apache Spark           & Data lake $\to$ Any   & Batch/stream & Large-scale transformation and enrichment \\
dbt                    & Data lake/DW internal & Batch    & SQL-based transformation, testing, lineage \\
Apache Atlas           & Any                   & Metadata  & Data governance, lineage graph, cataloguing \\
\bottomrule
\end{tabular}
\end{table}

A typical production ELT stack for a modern data-driven organisation
operates as follows. Source systems (RDBMS, SaaS applications, IoT
devices) emit data either via batch exports or real-time event streams.
Batch sources are ingested using Sqoop (for RDBMS-to-HDFS transfer) or
a managed connector platform (Airbyte, Fivetran) that handles
incremental extraction and schema change detection. Real-time sources
are ingested via \tool{Apache Kafka}, which acts as a durable, ordered
event bus with configurable retention---allowing consumers to replay
the full event history at any time. Kafka is examined in depth in
Chapter~\ref{ch:pipelines}.

Once raw data lands in the data lake (HDFS or cloud object storage),
transformation is managed by dbt (for SQL-based transformations in a
data warehouse engine) or Apache Spark (for large-scale Python/Scala
transformations over raw files). The transformation layer produces
\term{data marts}---curated, domain-specific views of the integrated
data designed for specific analytical workloads. A metadata and lineage
platform such as Apache Atlas automatically captures the transformation
history, maintaining the provenance graph required for regulatory
compliance.

\term{Schema evolution}---the challenge of handling changes in source
schemas without breaking downstream consumers---is managed by a
\term{schema registry}. Confluent Schema Registry (for Kafka-based
pipelines) maintains versioned schemas for each topic, allowing
producers and consumers to negotiate schema compatibility at runtime.
When a source system adds a new field, the registry enforces
\emph{backward compatibility}: consumers using the old schema can still
read new records by treating the new field as null; producers using the
new schema can still write records that old consumers can parse.

%%====================================================================
\section{Case Study: US Healthcare Interoperability and HL7 FHIR}
\label{sec:healthcare-case}
%%====================================================================

\begin{casestudy}{Healthcare Interoperability and HL7 FHIR --- Integration at National Scale}

\textbf{Background.}
The United States healthcare system is one of the world's most data-rich
and most data-fragmented. A typical US hospital uses between 20 and 50
different software systems: an electronic health record (EHR) platform,
a laboratory information system, a radiology picture-archiving system,
a pharmacy management system, a billing and insurance system, a
clinical decision support system, and dozens of others. These systems
were built by different vendors, at different times, using different
data models, and have historically been unable to exchange patient data
without significant manual effort.

\smallskip
\textbf{The integration cost.}
A 2019 analysis by the Center for Health Information and Decision
Systems (University of Maryland) estimated that integration failures
cost US hospitals \textbf{\$8.3 billion annually} in duplicated tests,
delayed care, medication errors attributable to incomplete records, and
manual reconciliation labour. A patient transferred between two hospitals
using different EHR vendors may have their medication list, allergy
history, and recent lab results effectively unavailable to the receiving
physician---an information gap that directly affects clinical decisions.

\smallskip
\textbf{The technical challenge as a heterogeneity problem.}
Healthcare data integration exhibits all four dimensions of
heterogeneity simultaneously. \emph{Structural}: patient identifiers
vary by hospital (medical record number, patient account number,
national identifier). \emph{Semantic}: the same diagnosis code may mean
different things under different coding systems (ICD-9 vs.\ ICD-10 vs.\
SNOMED CT). \emph{Syntactic}: dates, dosages, and units of measurement
are encoded differently across vendors. \emph{Granularity}: laboratory
reference ranges vary by testing equipment and laboratory protocol.

\smallskip
\textbf{The HL7 FHIR standard.}
Health Level 7 (HL7) is the standards organisation that developed the
Fast Healthcare Interoperability Resources (FHIR) standard to address
these challenges. FHIR defines a RESTful API and a set of standardised
data resources (Patient, Observation, Medication, DiagnosticReport,
etc.) that any EHR system can implement. By exposing patient data
through a FHIR API, a hospital allows any authorised consumer---a mobile
app, a referral system, a population health analytics platform---to
query standardised patient records without knowing the internal database
schema of the EHR.

FHIR resources are represented in JSON or XML. A standardised patient
resource looks like:

\begin{lstlisting}[style=pseudocode, caption={A partial HL7 FHIR Patient resource (JSON)}]
{
  "resourceType": "Patient",
  "id":          "pat-001",
  "name": [{"family": "Smith", "given": ["John"]}],
  "birthDate":   "1978-04-22",
  "gender":      "male",
  "identifier": [
    {"system": "urn:oid:2.16.840.1.113883.4.1",
     "value": "123-45-6789"}
  ],
  "address": [
    {"line": ["123 Main Street"],
     "city": "Boston", "state": "MA",
     "postalCode": "02101"}
  ]
}
\end{lstlisting}

\textbf{Integration architecture.}
A national health information exchange built on FHIR operates as
follows. Each participating hospital implements a FHIR server that
exposes its EHR data through the standard API. A health information
exchange (HIE) platform federates queries across multiple FHIR servers,
performing entity resolution (matching the same patient across hospitals
using demographic attributes) and aggregating standardised data into a
unified longitudinal patient record. The transformation from each
hospital's internal schema to FHIR is handled by a vendor-specific
FHIR adapter---essentially an ETL pipeline that runs inside or alongside
the EHR system.

\smallskip
\textbf{Data quality and entity resolution challenges.}
Even with FHIR standardisation, entity resolution across hospitals
remains the hardest technical challenge. The US has no universal patient
identifier (comparable to a national insurance number in many other
countries); matching patients across hospitals relies on probabilistic
matching on name, date of birth, address, and partial social security
number---exactly the similarity-function approach described in
Section~\ref{sec:er}. Typical production matching systems achieve 95--98\%
recall at 98--99\% precision, meaning that 1--5\% of true patient matches
are missed and 1--2\% of declared matches are incorrect---a significant
error rate in clinical contexts.

\smallskip
\textbf{Lessons for data engineering.}
The healthcare interoperability challenge illustrates three principles
that apply across all large-scale data integration projects. First,
technical standards alone (FHIR) are necessary but not sufficient;
data quality, entity resolution, and governance must be addressed
alongside standardisation. Second, the cost of integration failure is
not abstract---it manifests in clinical outcomes, regulatory penalties,
and operational waste. Third, integrating data from autonomous,
heterogeneous systems at scale requires exactly the distributed tools---
Kafka, Spark, schema registries, and lineage tracking---that constitute
modern data engineering practice.
\end{casestudy}

%%====================================================================
\section*{Chapter Summary}
\label{sec:ch2-summary}
%%====================================================================

\begin{itemize}
  \item \textbf{The data silo problem} is the primary barrier to
        enterprise analytics: data spread across independent, autonomous
        systems with incompatible schemas, identifiers, and formats
        cannot be jointly queried without integration.

  \item \textbf{Data integration} provides a unified query interface
        over heterogeneous, autonomous sources via three core operations:
        schema mapping (defining field correspondences), data exchange
        (materialising transformed data), and entity resolution
        (linking records representing the same real-world entity).

  \item \textbf{Four dimensions of heterogeneity} must all be addressed:
        structural (different schemas), semantic (different meanings),
        syntactic (different encodings), and scale/granularity (different
        units and resolutions). Structural and syntactic failures are
        loud; semantic and granularity failures are silent and
        dangerous.

  \item \textbf{Entity resolution} identifies matching records across
        sources with no common key. Blocking reduces the $O(n^2)$
        comparison problem; similarity functions (Jaccard, edit distance)
        score candidate pairs; a composite threshold determines matches.

  \item \textbf{Data quality} has six dimensions: accuracy, completeness,
        consistency, timeliness, validity, and uniqueness. Missing values
        may be MCAR, MAR, or MNAR; the appropriate imputation strategy
        depends on the missingness mechanism.

  \item \textbf{Data provenance} records the origin and transformation
        history of all data artifacts. It is required by GDPR and HIPAA
        and is operationally critical for debugging data pipelines.
        OpenLineage, Apache Atlas, and dbt provide lineage tracking.

  \item \textbf{ETL} transforms data before loading (high quality at
        the target, rigid schema). \textbf{ELT} loads raw data first
        and transforms on demand (flexible, supports schema evolution,
        requires a separate transformation layer such as dbt).

  \item The \textbf{modern integration stack} combines Kafka
        (streaming ingestion), Sqoop/Airbyte (batch ingestion),
        Spark/dbt (transformation), and Atlas (lineage). Schema
        registries manage schema evolution across producers and consumers.
\end{itemize}

%%====================================================================
\section*{Exercises}
\label{sec:ch2-exercises}
%%====================================================================

\exercisesection{Short Questions}

\sq Define data integration. What three operations does it require, and
what does each operation accomplish?

\sq Explain the difference between structural and semantic
heterogeneity. Give one example of each from the domain of retail
e-commerce.

\sq A source system stores dates as the string \texttt{"15/03/2024"}
and a target system expects dates as \texttt{DATE} values in ISO 8601
format (\texttt{"2024-03-15"}). Which dimension of heterogeneity does
this represent? What operation resolves it?

\sq What is entity resolution? Why is the naive $O(n^2)$ approach
infeasible at big data scale, and how does blocking address this?

\sq Define the six dimensions of data quality and give one concrete
example for each from a financial services context.

\sq What is the difference between MCAR, MAR, and MNAR missingness?
Why does the distinction matter when choosing an imputation strategy?

\sq What is data provenance? Name two regulatory frameworks that
require it and explain what each framework demands.

\sq Compare ETL and ELT: (a) in which order are the transform and load
steps applied, (b) where is the data quality enforcement point, and (c)
which is better suited for schema evolution?

\sq True or false, with one-sentence justification: ``The ELT paradigm
eliminates the need for data transformation.''

\sq What is a schema registry? What problem does it solve in a
Kafka-based integration pipeline?

\sq A retail company finds that 15\% of rows in its integrated customer
table have \texttt{NULL} in the \texttt{email} column. List three
possible reasons for this missingness and explain how you would
distinguish between them.

\sq What is ``schema drift''? Give a concrete example of how schema
drift can cause a downstream machine learning model to fail silently.

\sq What is the GAV (global-as-view) approach to schema mapping? How
does it differ from LAV (local-as-view), and what is the practical
trade-off between them?

\sq A data engineer says: ``Our ETL pipeline runs nightly; by morning
we have yesterday's data.'' Identify the data quality dimension being
compromised and explain what architecture change would address it.

\bigskip
\exercisesection{Analytical Problems}

\ap \textbf{Entity Resolution Complexity Analysis.}
A hospital system wants to resolve patient identities across two EHR
databases: Hospital A has 2.8 million patient records; Hospital B has
1.4 million patient records. A blocking strategy based on birth-year
and first three characters of the last name produces blocks of average
size 350 pairs (350 A-records times the B-records in the same block).

\begin{enumerate}[label=(\alph*)]
  \item Calculate the number of candidate pairs under naive $O(n^2)$
        comparison (all A against all B). At a rate of 5 million
        comparisons per second, how long would naive comparison take?
  \item Estimate the total number of candidate pairs after blocking,
        assuming each blocking key creates blocks of average size 350.
        (Hint: how many blocks are there if each A-record falls in
        exactly one block?)
  \item What is the speedup factor from blocking?
  \item A composite similarity score requires 12 floating-point
        operations per pair comparison. Estimate the total computation
        cost (in FLOPs) for the post-blocking comparisons, and determine
        whether it is feasible on a single server with 100 GFLOPS of
        sustained throughput.
\end{enumerate}

\ap \textbf{ETL vs.\ ELT Storage and Latency Trade-off.}
A financial services company ingests 5 TB of raw trading data per day
from 40 source systems. In the current ETL architecture, data is
transformed and compressed before loading, resulting in a 4:1
compression ratio at the warehouse. The company is evaluating
migrating to an ELT architecture with raw data stored in S3 at
\$0.023/GB/month, and transformation via Spark on demand.

\begin{enumerate}[label=(\alph*)]
  \item Calculate the current warehouse storage cost after 12 months
        (assume warehouse storage at \$0.20/GB/month, and that the
        compression ratio applies to all stored data). Compare with the
        raw data storage cost in S3 over the same period.
  \item The ETL pipeline requires 4 hours of processing before data is
        available for queries. The ELT approach makes raw data available
        immediately but requires 45 minutes of Spark transformation
        before a specific analytical view is ready. For a use case that
        needs data refreshed every 6 hours, which approach provides
        better effective freshness?
  \item In the ELT system, 12 months of raw data are retained in S3.
        A new analytical requirement arrives, requiring a transformation
        that was not anticipated. Estimate the Spark compute cost to
        retroactively apply this transformation to 12 months of 5 TB/day
        data at \$0.10 per Spark CPU-hour (assume 10 TB/hour processing
        throughput per 10 CPUs).
\end{enumerate}

\ap \textbf{Data Quality Profiling.}
An e-commerce company's data engineering team profiles a 50-million-row
product catalogue table and finds the following:
\begin{itemize}
  \item 12\% of rows have \texttt{NULL} in \texttt{product\_weight}.
  \item 3.2\% of rows have \texttt{price} $= 0$.
  \item 0.8\% of rows have \texttt{price} $> \$10{,}000$ (the 99.9th
        percentile is \$950).
  \item 4\% of rows appear to be duplicates (same product, same vendor,
        inserted on different dates).
  \item 7\% of rows have \texttt{category\_id} values that do not exist
        in the \texttt{categories} reference table.
\end{itemize}
For each issue: (a) identify the data quality dimension it violates,
(b) classify it as noise, missing data, inconsistency, or validity
violation, and (c) propose a specific remediation and the operational
procedure to apply it without breaking downstream pipelines.

\ap \textbf{Similarity Scoring for Entity Resolution.}
Two records from different source systems are candidates for entity
resolution. Using the composite similarity formula with weights
$w_\text{name} = 0.5$, $w_\text{dob} = 0.3$, $w_\text{city} = 0.2$:

\medskip
\noindent
\begin{tabular}{lll}
\toprule
\textbf{Attribute} & \textbf{Source A} & \textbf{Source B} \\
\midrule
Name      & \texttt{Johnson Michael R.} & \texttt{Michael R Johnson} \\
Date of birth & \texttt{1985-07-14} & \texttt{07/14/1985} \\
City      & \texttt{Los Angeles} & \texttt{LA} \\
\bottomrule
\end{tabular}

\medskip
\noindent
(a) Parse both date representations and confirm they are the same date
($\sigma_\text{dob} = 1.0$). (b) Compute the Jaccard similarity on
name tokens (split on space and punctuation); let
$\sigma_\text{name} = J(T_A, T_B)$. (c) Assign
$\sigma_\text{city} = 0.4$ (``Los Angeles'' and ``LA'' are the same
city but share no token; domain knowledge is required). (d) Compute
$\sigma_\text{composite}$ and determine whether these records are a
match at thresholds $\theta = 0.7$ and $\theta = 0.8$. (e) Discuss
whether the threshold choice here is an engineering or a business
decision.

\bigskip
\exercisesection{Design Problems}

\dpr \textbf{Integration Architecture for a Global Logistics Company.}
A global freight logistics company operates in 60 countries. Each
country office uses a locally deployed ERP system (Oracle in North
America, SAP in Europe, a custom-built system in Southeast Asia). The
company needs a unified analytics platform that answers questions such
as: ``Which shipment routes had the highest delay rates last quarter?''
and ``What is the total revenue per customer across all regions?''

The three ERP systems use different customer identifiers, different
currency representations (some systems store amounts in local currency
only; others in USD), different status code vocabularies for shipment
states, and different date formats. Real-time track-and-trace data
(GPS position of shipments) arrives from 8,000 vehicles at 30-second
intervals.

Design the integration architecture. Your design must specify:
\begin{enumerate}[label=(\alph*)]
  \item The integration paradigm (ETL or ELT) and justification.
  \item How each of the four heterogeneity dimensions is addressed.
  \item The entity resolution strategy for customers (who may appear
        in multiple regional ERP systems with different identifiers).
  \item How the real-time GPS stream is ingested and at what latency
        it becomes available for analytics.
  \item How data lineage and provenance are maintained across the pipeline.
  \item Which specific tools you would use at each stage and why.
\end{enumerate}

\dpr \textbf{Healthcare Data Quality Framework.}
A national health analytics agency receives patient-level data files
from 200 hospitals every month. The files arrive as CSV exports from
each hospital's EHR system. The agency uses this data to compute
national statistics on chronic disease prevalence, hospital readmission
rates, and treatment outcome metrics.

A preliminary audit reveals: (1) 18 different date formats across the
200 hospitals; (2) ICD-10 codes mixed with ICD-9 codes in the same
diagnosis field; (3) patient ages ranging from --3 to 247 (clearly
erroneous); (4) 22\% of records missing at least one of: gender,
ethnicity, or primary diagnosis; and (5) duplicate patient records at
rates between 0.5\% and 8\% depending on the hospital.

Design a data quality pipeline that the agency can apply to incoming
files before they are used in national statistics. Your design must:
\begin{enumerate}[label=(\alph*)]
  \item Specify a quality gate for each of the five issues identified,
        including the rule, the action for records that fail (reject,
        flag, impute), and the threshold at which an entire hospital's
        submission is quarantined.
  \item Describe how MCAR vs.\ MAR vs.\ MNAR missingness would be
        diagnosed for the 22\% missing records, and what imputation
        strategy is appropriate for each case.
  \item Explain how you would maintain provenance to satisfy HIPAA
        audit requirements while working with de-identified patient data.
  \item Propose a data quality SLA (e.g., ``we require $>$98\%
        completeness on primary diagnosis before publishing national
        statistics'') and explain how it would be monitored in
        production.
\end{enumerate}

\bigskip
\exercisesection{Programming Exercises}

\pe \textbf{Entity Resolution with Jaccard Similarity (Python).}
Implement the Jaccard similarity function for string tokenisation and
apply it to resolve company names across two source datasets.

\begin{lstlisting}[style=pyspark, caption={PE2.1: Company name entity resolution}]
import re

source_a = [
    {"id": "A001", "name": "Google LLC",       "city": "Mountain View"},
    {"id": "A002", "name": "Apple Inc.",        "city": "Cupertino"},
    {"id": "A003", "name": "Microsoft Corp",    "city": "Redmond"},
    {"id": "A004", "name": "Meta Platforms Inc","city": "Menlo Park"},
    {"id": "A005", "name": "Amazon.com Inc",    "city": "Seattle"},
]

source_b = [
    {"id": "B001", "name": "GOOGLE LLC",          "city": "Mountain View"},
    {"id": "B002", "name": "Apple Incorporated",  "city": "Cupertino"},
    {"id": "B003", "name": "Microsoft Corporation","city": "Redmond"},
    {"id": "B004", "name": "Facebook (Meta)",      "city": "Menlo Park"},
    {"id": "B005", "name": "Amazon Web Services",  "city": "Seattle"},
]

def tokenise(text):
    """Lowercase, remove punctuation, split into tokens."""
    text = re.sub(r"[.,\-()']", " ", text.lower())
    return set(text.split())

def jaccard(tokens_a, tokens_b):
    """
    Return Jaccard similarity between two token sets.
    J(A, B) = |A intersection B| / |A union B|
    Return 0.0 if the union is empty.
    """
    # TODO: implement
    pass

def blocking_key(record):
    """
    Return a blocking key: first token of name + first 3 chars of city.
    Records only compared within the same blocking key.
    """
    first_token = tokenise(record["name"])
    city_prefix = record["city"][:3].lower()
    # TODO: return key string
    pass

def resolve_entities(source_a, source_b, threshold=0.5):
    """
    Return list of (a_id, b_id, score) tuples where score >= threshold.
    Use blocking to reduce comparisons.
    """
    matches = []
    # TODO:
    # 1. Build index of source_b records by blocking key
    # 2. For each record in source_a, compute its blocking key
    # 3. Compare only against source_b records with the same key
    # 4. Compute Jaccard on name tokens; append match if >= threshold
    return matches

matches = resolve_entities(source_a, source_b, threshold=0.4)
print("Resolved entity pairs:")
for a_id, b_id, score in matches:
    a_rec = next(r for r in source_a if r["id"] == a_id)
    b_rec = next(r for r in source_b if r["id"] == b_id)
    print(f"  {a_id} ({a_rec['name']}) <-> {b_id} ({b_rec['name']})  "
          f"score={score:.3f}")
\end{lstlisting}

\pe \textbf{Data Quality Audit (Python).}
Implement a data quality profiler that checks a list of records against
defined quality rules and produces a summary report.

\begin{lstlisting}[style=pyspark, caption={PE2.2: Data quality profiler}]
import datetime

records = [
    {"id": "R001", "name": "Alice Chen",
     "dob": "1990-05-14", "email": "alice@example.com", "amount": 1250.00},
    {"id": "R002", "name": "Bob Smith",
     "dob": "1985-13-01",  # invalid month
     "email": None, "amount": 0.0},
    {"id": "R003", "name": "",            # empty name
     "dob": "2035-01-01",  # future date
     "email": "carol@example.com", "amount": -50.0},
    {"id": "R004", "name": "Diana Kumar",
     "dob": "1978-03-22", "email": "diana@example.com", "amount": 980.50},
    {"id": "R005", "name": "Eve Johnson",
     "dob": "1978-03-22",  # duplicate dob+name (same as R004 roughly)
     "email": None, "amount": 99999.0},
]

def check_completeness(record, required_fields):
    """Return list of field names that are None or empty string."""
    missing = []
    for f in required_fields:
        v = record.get(f)
        if v is None or v == "":
            missing.append(f)
    return missing

def check_dob_validity(dob_str):
    """
    Return True if dob_str is a valid date in ISO format (YYYY-MM-DD)
    and not in the future.
    """
    try:
        dob = datetime.date.fromisoformat(dob_str)
        return dob <= datetime.date.today()
    except (ValueError, TypeError):
        return False

def check_amount_validity(amount):
    """Return True if amount is a positive number."""
    try:
        return float(amount) > 0
    except (ValueError, TypeError):
        return False

def profile_records(records):
    """
    Run all quality checks and print a summary report.
    Dimensions checked: Completeness, Validity, Accuracy (range).
    """
    required = ["name", "dob", "email", "amount"]
    issues = []

    for rec in records:
        rec_issues = {"id": rec["id"], "problems": []}

        # Completeness
        missing = check_completeness(rec, required)
        if missing:
            rec_issues["problems"].append(
                f"Missing fields: {missing}")

        # Validity: dob
        if rec.get("dob") and not check_dob_validity(rec["dob"]):
            rec_issues["problems"].append(
                f"Invalid or future dob: {rec['dob']}")

        # Validity: amount
        if rec.get("amount") is not None:
            if not check_amount_validity(rec["amount"]):
                rec_issues["problems"].append(
                    f"Non-positive amount: {rec['amount']}")

        if rec_issues["problems"]:
            issues.append(rec_issues)

    # Summary
    total = len(records)
    clean = total - len(issues)
    print(f"Quality Report: {clean}/{total} records passed all checks\n")
    for issue in issues:
        print(f"  Record {issue['id']}:")
        for p in issue["problems"]:
            print(f"    - {p}")
    return issues

profile_records(records)
\end{lstlisting}

\pe \textbf{ETL Pipeline Simulation (Python).}
Simulate a three-stage ETL pipeline that extracts records from two
heterogeneous sources, applies transformations to resolve structural
and syntactic heterogeneity, and loads unified records into a target
schema.

\begin{lstlisting}[style=pyspark, caption={PE2.3: ETL pipeline simulation}]
import datetime

# Source A: CRM system (European date format, price in EUR)
source_a = [
    {"contact_id": "C001", "fname": "Emma",  "lname": "Weber",
     "birth": "22.09.1988", "eur_value": 1400.0},
    {"contact_id": "C002", "fname": "Luca",  "lname": "Rossi",
     "birth": "05.03.1975", "eur_value": 2200.0},
]

# Source B: ERP system (US date format, price in USD, full name column)
source_b = [
    {"customer_no": "E045", "full_name": "Chen, Wei",
     "date_of_birth": "07/18/1992", "usd_value": 3100.0},
    {"customer_no": "E046", "full_name": "Park, Ji-Young",
     "date_of_birth": "12/04/1983", "usd_value": 870.0},
]

EUR_TO_USD = 1.09  # approximate exchange rate

# Target schema:
# { "id", "first_name", "last_name", "dob" (ISO), "value_usd" }

def transform_source_a(record):
    """
    Map CRM record to target schema.
    - contact_id -> id
    - fname/lname -> first_name/last_name
    - birth (DD.MM.YYYY) -> dob (YYYY-MM-DD)
    - eur_value * EUR_TO_USD -> value_usd
    """
    # TODO: implement transformation
    pass

def transform_source_b(record):
    """
    Map ERP record to target schema.
    - customer_no -> id
    - full_name ("Last, First") -> first_name, last_name
    - date_of_birth (MM/DD/YYYY) -> dob (YYYY-MM-DD)
    - usd_value -> value_usd (already USD)
    """
    # TODO: implement transformation
    pass

def run_etl(source_a, source_b):
    """Extract, transform both sources, and load into unified target."""
    target = []

    # Transform source A
    for rec in source_a:
        transformed = transform_source_a(rec)
        if transformed:
            target.append(transformed)

    # Transform source B
    for rec in source_b:
        transformed = transform_source_b(rec)
        if transformed:
            target.append(transformed)

    return target

unified = run_etl(source_a, source_b)
print("Unified target records:")
for r in unified:
    print(f"  {r}")
\end{lstlisting}
